\documentclass[a4paper, 12pt, english]{article}

\usepackage[utf8]{inputenc}
\usepackage[T1]{fontenc}
\usepackage{babel}

% Use Garamond font
\usepackage{ebgaramond}
\usepackage{float}
\usepackage{amsmath,amssymb}
\usepackage{graphicx}
\usepackage{subfig}
\usepackage[colorinlistoftodos]{todonotes}
\usepackage{matlab-prettifier}
\usepackage{indentfirst}
\usepackage{verbatim}
\usepackage{textcomp}
\usepackage{gensymb}
\usepackage{longtable}

\usepackage{relsize}

\usepackage{lipsum}
\usepackage{xcolor}
\usepackage{xparse}
\NewDocumentCommand{\myrule}{O{1pt} O{2pt} O{black}}{%
  \par\nobreak % don't break a page here
  \kern\the\prevdepth % don't take into account the depth of the preceding line
  \kern#2 % space before the rule
  {\color{#3}\hrule height #1 width\hsize} % the rule
  \kern#2 % space after the rule
  \nointerlineskip % no additional space after the rule
}
\usepackage[section]{placeins}

\usepackage{booktabs}
\usepackage{colortbl}%
   \newcommand{\myrowcolour}{\rowcolor[gray]{0.925}}
   
\usepackage[obeyspaces]{url}
\usepackage{etoolbox}
\usepackage[colorlinks=true, linkcolor=black]{hyperref}

\usepackage{geometry}
\geometry{
    paper=a4paper, % Change to letterpaper for US letter
    inner=1.8cm, % Inner margin
    outer=1.8cm, % Outer margin
    bindingoffset=.5cm, % Binding offset
    top=1.8cm, % Top margin
    bottom=1.8cm, % Bottom margin
    %showframe, % Uncomment to show how the type block is set on the page
}
%*******************************************************************************%
%************************************START**************************************%
%*******************************************************************************%
\begin{document}

\begin{titlepage}
\begin{center}
\par
\vfill
\vspace{500pt}
\includegraphics[width=0.65\textwidth]{figs/bath22.png}

\vspace{70pt}
\textbf{\LARGE  MESO FIBRE PARTICLE SIZE CHARACTERISATION}\\
\vspace{10pt}
\textbf{\large System Requirements}\\
\vspace{30pt}
\begin{table}[H]
  \centering
  \renewcommand{\arraystretch}{1.2}
  \begin{tabular}{@{}llllp{6.5cm}@{}}
    \toprule
    \textbf{Version} & \textbf{Date} & \textbf{Author} & \textbf{Reviewed By} & \textbf{Change Summary} \\
    \midrule
    0.1 & 2026-02-09 & Van Kaliafetis & Joshua Poole & Initial requirements proposed. \\
    0.2 & 2026-03-06 & Van Kaliafetis & Ori Sela & Restructured requirements and added assumptions. \\
    1.0 & 2026-03-11 & Van Kaliafetis & Joshua Poole & Finalised requirements and formatting. \\
    \bottomrule
  \end{tabular}
\end{table}

\par
\vfill
\vspace{275pt}

\par
\vfill
\end{center}

\end{titlepage}

\newpage

\tableofcontents
\newpage

\setcounter{tocdepth}{2}

\section{Scope}
\subsection{System Purpose}
The system shall provide laboratory-scale characterisation of dry recycled textile fibres intended for use as partial replacement of virgin wood fibres in Medium Density Fibreboard (MDF) manufacture.

The system shall:

\begin{itemize}
    \item Employ a modular architecture to enable straightforward integration into existing manufacturing workflows and reduce barriers to commercial adoption.
    \item Quantify fibre geometric properties (e.g. length, width proxy, aspect ratio).
    \item Provide fibre size distribution suitable for assessing blend suitability and process risk.pilot-scale decision-making
    \item The system shall minimise capital and operational costs to ensure economic feasibility for deployment within MDF manufacturing facilities.
\end{itemize}

The system is not intended to:

\begin{itemize}
    \item Perform analysis on virgin feedstock.
    \item Produce MDF boards.
    \item Replace full board-level mechanical testing.
    \item Alter the properties of a sample
    \item Perform any separation or blending of samples.
\end{itemize}


\section{Assumptions}

The following assumptions underpin the requirements:

\begin{itemize}
    \item \textbf{ASSUMP-01} - Fibre samples are delivered dry and free-flowing, having already been processed through a textile milling system.
    \item \textbf{ASSUMP-02 }- Samples may contain heterogeneous mixtures of fibre types, colours, and fineness typical of post-consumer textile waste.
    \item \textbf{ASSUMP-03} - The primary measurement objective is geometric characterisation (size and morphology).
    \item \textbf{ASSUMP-04} - The system will operate in a standard indoor industrial environment with stable ambient temperature and non-condensing humidity.
    \item \textbf{ASSUMP-05} - The system will be operated by trained industrial users (e.g. engineers or researchers), not by untrained members of the public.
    \item \textbf{ASSUMP-06} - Statistical relevance is achieved through sampling and aggregation of many fibres rather than complete measurement of every fibre in a batch.
    \item \textbf{ASSUMP-07} - Image analysis may use deterministic algorithms and/or AI-based approaches, but final outputs must remain interpretable and traceable.
    \item \textbf{ASSUMP-08} – Typical plants process approximately 90 kg/min virgin wood fibre input. Substituting approximately 30\% (Fab Materials estimate) of this stream leads to a system throughput of approximately 27 kg/min.
\end{itemize}

\section{System Requirements}

\subsection{Definitions}

\begin{itemize}
    \item \textbf{Batch:} A total mass of textile that must be characterized.
    \item \textbf{Sample:} Taken from a batch. Samples are analyzed for size characterization.
\end{itemize}

\subsection{High Level System Requirements}

\begin{itemize}
    \item \textbf{}\textbf{SYS-FUNC-01}- The system shall accept and process non-homogeneous (mixed composition) textile batches. 
    \item \textbf{SYS-FUNC-02}- The system shall obtain representative samples from textile batches. 
    \item \textbf{SYS-FUNC-03}- The system shall make fibre size distribution data for analysed samples available to the user. 
    \item \textbf{SYS-FUNC-04}- The system shall make modelled board output characteristics (using the current sample) available to the user. 
    \item \textbf{SYS-FUNC-05}- The system shall produce consistent results when analyzing equivalent samples, with a coefficient of variation <5\%. 
    \item \textbf{SYS-FUNC-06}- The system shall maintain traceability between batches, samples, and analysis results. 
    \item \textbf{SYS-FUNC-07}- The system shall characterise fibre samples at a throughput of at least 27 kg/min. 
    \item \textbf{SYS-FUNC-08}- The system shall be constructed with components with a total cost not exceeding £5,000 [5k is variable]. 
\end{itemize}

\subsection{Measurement \& Performance Requirements}

\begin{itemize}
    \item \textbf{SYS-PERF-01}- The system shall measure individual fibres within a length range of approximately 1 mm to 10 mm and a width range of 0.05 mm to 0.2 mm.

    \item \textbf{SYS-PERF-02}- The system shall estimate fibre length with a systematic error of less than $\pm$5\% for fibres within the nominal measurement range.

    \item \textbf{SYS-PERF-03}- The system shall estimate fibre width (or equivalent diameter proxy) with a systematic error not exceeding $\pm$15\%.

    \item \textbf{SYS-PERF-04}- For repeated measurements of the same sample, the coefficient of variation (CV) of mean fibre dimensions shall not exceed 5\%.

    \item \textbf{SYS-PERF-06}- The system shall maintain calibration drift within $\pm$5\% over a standard laboratory operating week without recalibration.
\end{itemize}

\subsection{Operational \& Throughput Requirements}

\begin{itemize}
    \item \textbf{SYS-OPER-01}- The system shall operate for a full industrial working period without unscheduled maintenance.
    \item \textbf{SYS-OPER-02}- The system shall allow safe interruption and restart of measurements without data corruption.
\end{itemize}

\subsection{Data \& AI Requirements}

\begin{itemize}
    \item \textbf{SYS-DATA-01}- The system shall keep a cache of recent measurements to defend against interruption.
    \item \textbf{SYS-DATA-02}- The system shall protect stored data from unintended modification.
    \item \textbf{SYS-DATA-03}- If AI-based models are used, the system shall maintain model version traceability.
    \item \textbf{SYS-DATA-04}- The system shall provide user access control appropriate to industrial usage.
\end{itemize}

\subsection{Human--Machine Interface Requirements}

\begin{itemize}
    \item \textbf{SYS-HMI-01}- The system shall clearly indicate operational state (idle, running, fault, calibration).
    \item \textbf{SYS-HMI-02}- The system shall provide clear software error messages and recommended operator actions.
\end{itemize}

\subsection{Physical \& Environmental Requirements}

\begin{itemize}
    \item \textbf{SYS-PHYS-01}- The system shall operate in standard indoor industrial conditions.
    \item \textbf{SYS-PHYS-02}- The system shall have dimensions and mass compatible with typical industrial handling.
    \item \textbf{SYS-PHYS-03}- The system shall control or shield the imaging region from uncontrolled ambient light.
    \item \textbf{SYS-PHYS-04}- The system shall include provisions to manage fibre dust accumulation within the enclosure.
    \item \textbf{SYS-PHYS-05}- The system shall limit noise and vibration to acceptable levels (threshold to be specified) so imaging is not degraded.

\end{itemize}

\subsection{Safety Requirements}

\begin{itemize}
    \item \textbf{SYS-SAFE-01}- The system shall prevent operator access to hazardous moving parts during operation.
    \item \textbf{SYS-SAFE-02}- The system shall shield users from high-intensity illumination sources where applicable.
    \item \textbf{SYS-SAFE-03}- The system shall include an emergency stop function where moving mechanisms are present.
    \item \textbf{SYS-SAFE-04}- The system shall comply with applicable electrical and laboratory equipment safety standards.
    \item \textbf{SYS-SAFE-05}- The system shall control airborne fibre release at all stages via enclosure and/or extraction provisions.  
    \item \textbf{SYS-SAFE-06}- The system shall provide a documented PPE requirement and safe operating procedure 
\end{itemize}

\subsection{Reliability, Maintainability \& Lifecycle}

\begin{itemize}
    \item \textbf{SYS-RAM-01}- The system shall be designed for reliable operation under normal industrial usage.
    \item \textbf{SYS-RAM-02}- The system shall allow routine cleaning and maintenance using standard tools.
    \item \textbf{SYS-RAM-03}- The system shall log operational and diagnostic information for fault investigation.
    \item \textbf{SYS-RAM-04}- The system shall allow recalibration using defined procedures.
    \item \textbf{SYS-RAM-05}- The system shall be designed to allow future upgrade of analysis algorithms without full hardware redesign.
    \item \textbf{SYS-RAM-06}- The system shall detect and flag conditions that invalidate results (e.g., calibration drift, poor illumination, excessive overlap)  
\end{itemize}

\end{document}
